\documentclass[11pt,a4paper]{article}

% --- packages -----------------------------------------------------------------
\usepackage[utf8]{inputenc}
\usepackage[T1]{fontenc}
\usepackage{microtype}
\usepackage{lmodern}
\usepackage[margin=1in]{geometry}
\usepackage{graphicx}
\usepackage{booktabs}
\usepackage{hyperref}
\usepackage{xcolor}
\usepackage{caption}
\usepackage{enumitem}
\usepackage[numbers,square]{natbib}

\hypersetup{
  colorlinks=true,
  linkcolor=blue!50!black,
  citecolor=blue!40!black,
  urlcolor=blue!50!black,
}

\captionsetup{font=small, labelfont=bf, skip=4pt}
\setlength{\parskip}{0.4em}
\setlength{\parindent}{0pt}

\newcommand{\graphicsdir}{assets}

% --- title block --------------------------------------------------------------
\title{\textbf{Fuzuli v0.1} \\[0.3em]
  \large A from-scratch Latinized Ottoman Turkish language model}

\author{
  Fatih Burak Karagöz\thanks{Contact: \texttt{fatih@cdli.ai}} \\[0.2em]
  \normalsize\textsc{CDLI}\thanks{CDLI is an independent laboratory for AI
    products and systems. \texttt{cdli.ai}.} \\
  \normalsize Istanbul, Turkey
}

\date{May 2026}

\begin{document}
\maketitle

\begin{abstract}
\noindent
\textbf{Fuzuli v0.1} is an 84-million-parameter decoder-only language model
trained from scratch on Latinized post-1500 Ottoman Turkish. The training
corpus consists of 24,253 OCR-cleaned pages (70.1 million characters,
${\approx}\,8$ million unique tokens under a custom 24,000-token BPE
vocabulary) drawn from 59 distinct historical documents — among them the full
available \emph{Seyahatnâme} of Evliya Çelebi, the complete 25-volume run of
the late-Ottoman periodicals \emph{Sebilürreşad} and \emph{Sırat-ı Müstakim},
and 25 individual classical and late-classical works ranging from Necati
Bey's \emph{Dîvân} (1500) to \emph{Üç Devirde Gördüklerim} (1866). No
foreign weights were inherited; no model was fine-tuned. The tokenizer, the
training data, and the model weights are all produced end-to-end inside the
project. This first version is released as a research artifact, not a
production tool. It demonstrates that a small from-scratch Ottoman language
model is feasible at the current data scale, and it documents an
\emph{autoresearch} methodology --- a deterministic factory paired with an
LLM-based research agent under strict eval-policy control --- that produced
six baseline-improving recipes out of 356 candidate experiments in the
model's first 24 hours of life.
\end{abstract}

% =============================================================================
\section{Why Ottoman Turkish, why from scratch}
% =============================================================================

Ottoman Turkish is not a single language, it is six centuries of overlapping
registers --- classical \emph{dîvân} poetry, \emph{seyahatnâme} travel prose,
\emph{tezkire} biographical compilations, late-Ottoman periodical journalism,
\emph{risâle} religious treatises --- written across a script transition
(Arabic-letter to Latin-letter, formalized in 1928) that broke continuity for
most modern computational tools. The corpus that does exist in
machine-readable Latinized form is small, scattered across academic
websites, and rarely curated for language modeling.

Fine-tuning a multilingual model is the easy path; it is also the path that
bakes in modern Turkish bias and imports tokenization choices made for
languages that share almost no morphology with Ottoman. The decision here is
the harder one: own the tokenizer, own the data pipeline, own every weight
from the first matmul forward. The cost is a smaller model trained on less
data; the benefit is an artifact that is unambiguously \emph{of} the corpus
it was trained on. This stance follows a wider tradition of from-scratch
low-resource language models --- BERTurk~\citep{schweter2020berturk}, the
OSCAR-trained European national models~\citep{suarez2019oscar}, and
curated-corpus small-LM training in the spirit of
TinyStories~\citep{eldan2023tinystories} --- which have repeatedly
demonstrated that a smaller, language-faithful model trained on a curated
in-domain corpus outperforms a larger multilingual model fine-tuned on the
same downstream task.

\subsection{A note on prior work}

Fuzuli continues a line of work I have contributed to on the computational
treatment of Ottoman Turkish. In \emph{Towards a Clean Text Corpus for
Ottoman Turkish}~\citep{karagoz2024towards}, my co-authors and I constructed
an initial cleaned Ottoman corpus and used it for \textbf{continual
pre-training} of BERTurk~\citep{schweter2020berturk}, demonstrating that
even modest amounts of in-domain text can adapt a Modern Turkish encoder to
a historical variant for downstream tasks such as named entity recognition.
The cleaning methodology developed there --- regex-based normalization,
handling of intertwined bidirectional Arabic/Latin script, private-use-area
character mapping, dehyphenation --- is the direct predecessor of
Section~\ref{sec:corpus}'s four-stage pipeline.

The broader infrastructure effort I co-author with the Boğaziçi NLP group,
\emph{Building Foundations for Natural Language Processing of Historical
Turkish: Resources and Models}~\citep{ozates2025foundations}, introduced the
\textbf{HisTR} dataset, the \textbf{OTA-BOUN} Universal Dependencies
treebank, and the \textbf{OTC} corpus, alongside fine-tuned task models
published under the \texttt{bucolin} organization on HuggingFace. That work
establishes concrete downstream baselines for historical Turkish --- NER F1
of \textbf{90.29} on HisTR, dependency-parsing LAS of \textbf{73.79} on
OTA-BOUN, POS F1 of \textbf{94.98} --- against which a from-scratch
decoder-only successor can be measured.

Fuzuli takes the next step in that programme. Rather than continually
pre-training or fine-tuning Modern Turkish base models, it trains a
decoder-only language model end-to-end on the historical variant alone ---
with its own tokenizer, its own frozen evaluation harness, and no inherited
weights. The progression is intentional: encoder fine-tuning answers
\emph{can the existing infrastructure be adapted?}, and the answer is
\emph{yes, partially}. From-scratch decoder pretraining answers a different
question --- \emph{what does the language look like to a model that has
never been told anything about Modern Turkish?} --- and the resulting
artifact is, by construction, free of any Modern-Turkish prior.

% =============================================================================
\section{The corpus}\label{sec:corpus}
% =============================================================================

The training data is the union of two HuggingFace datasets I previously
curated and published --- \texttt{fatihburakkaragoz/anadolu-ocr-corpus} (52
documents, OCR'd via DeepSeek-OCR-2 at 200--300 DPI) and
\texttt{fatihburakkaragoz/evliya-celebi-seyahatname-ocr} (7 books of Evliya
Çelebi's travelogue) --- combined and processed through a four-stage
cleaning pipeline:
\begin{enumerate}[leftmargin=*,nosep]
  \item Unicode NFC normalization, control-character stripping, whitespace
    collapse.
  \item Regex-based cleanup, length filters, modern-Turkish loanword ratio
    filter.
  \item MinHash near-duplicate removal at Jaccard $\geq 0.85$.
  \item Per-document train/heldout split
    (\texttt{sha256(source\_pdf)\,\%\,100\,<\,2} $\to$ heldout).
\end{enumerate}

After cleaning, the corpus contains \textbf{24{,}253 pages}, \textbf{70.1
million characters}, and --- once tokenized with the project's own 24k BPE
--- \textbf{7.97 million unique training tokens} plus 116{,}000 heldout
tokens.

\begin{figure}[t]
  \centering
  \includegraphics[width=0.65\textwidth]{\graphicsdir/corpus_era.png}
  \caption{Corpus composition by era (post-cleaning).}
  \label{fig:era}
\end{figure}

\begin{figure}[t]
  \centering
  \includegraphics[width=0.85\textwidth]{\graphicsdir/corpus_genre.png}
  \caption{Corpus composition by genre.}
  \label{fig:genre}
\end{figure}

Era is inferred from each document's date label or, when absent, from the
publication-year heuristic encoded in the importer (with a hard-coded
override for known periodicals whose filenames carry misleading print-run
dates). The actual distribution --- 59\% classical, 40\% late Ottoman,
1\% Tanzimat --- reflects the source material rather than a deliberate
sampling choice.

\begin{figure}[t]
  \centering
  \includegraphics[width=\textwidth]{\graphicsdir/corpus_top_documents.png}
  \caption{Top-20 documents in the post-cleaning corpus, by character count.}
  \label{fig:topdocs}
\end{figure}

\subsection{Tokenizer: why train your own}\label{sec:tokenizer}

\begin{figure}[t]
  \centering
  \includegraphics[width=0.85\textwidth]{\graphicsdir/tokenizer_compression.png}
  \caption{Tokenization efficiency on the Fuzuli training corpus (28.8\,M
    characters of post-cleaning Latinized Ottoman). Custom BPE compresses
    1.74$\times$ better than OpenAI's \texttt{cl100k\_base} on the same text.
    Note: \texttt{cl100k} is not trained on Ottoman; the comparison is
    intended only as a sanity-check on language-specific tokenizer benefit.}
  \label{fig:tokenizer}
\end{figure}

A custom Byte-Pair Encoding tokenizer~\citep{sennrich2016bpe} of vocabulary
24{,}000 was trained on the post-cleaning train split. On the same Ottoman
corpus the custom BPE compresses at \textbf{3.61 chars/token}, against
\textbf{2.07} for OpenAI's \texttt{cl100k\_base} and \textbf{2.38} for the
newer \texttt{o200k\_base} --- a \textbf{1.74$\boldsymbol\times$ advantage}
over a general-purpose tokenizer. The reason is straightforward: Ottoman
morphology contains long agglutinative chains
(\textit{-larımızdaki}, \textit{-mıştır}) and Persian/Arabic loan-stems
(\textit{hüsn-ü}, \textit{eser-i}) that general-purpose tokenizers
fragment into bytes, while a corpus-trained vocabulary captures them as
single subword pieces. The 24k vocabulary size is on the smaller end (LLaMA-2
uses 32k~\citep{touvron2023llama2}, GPT-4 uses 100k+); with only 8\,M
unique training tokens, a 24k vocabulary places average occurrence per
vocab entry at the lower edge of healthy BPE coverage ($\sim$330 instances
per piece).

% =============================================================================
\section{Architecture}\label{sec:architecture}
% =============================================================================

Fuzuli is a standard post-2023 dense decoder transformer. None of its
individual components are novel; what is intentional is the
\emph{combination} and the discipline of refusing any architectural choice
that requires a foreign-weight initializer.

\begin{figure}[t]
  \centering
  \includegraphics[width=\textwidth]{\graphicsdir/model_architecture.png}
  \caption{Fuzuli architecture: one decoder block (left) and the full 12-block
    stack (right). Standard post-2023 dense decoder transformer with
    RMSNorm, RoPE, GQA, and SwiGLU.}
  \label{fig:arch}
\end{figure}

\begin{table}[t]
  \centering
  \small
  \begin{tabular}{ll}
    \toprule
    \textbf{Component} & \textbf{Choice} \\
    \midrule
    Layers & 12 \\
    Hidden dim & 640 \\
    Attention heads (Q) & 8 \\
    KV heads (GQA) & 4 \\
    Head dim & 80 \\
    MLP & SwiGLU, mlp\_ratio 2.67 (hidden 1707) \\
    Positional embedding & RoPE, $\theta = 10{,}000$ \\
    Normalization & RMSNorm, no bias \\
    Embedding tying & False (untied LM head) \\
    Tokenizer & BPE, vocab 24{,}000, byte-fallback \\
    Total parameters & $\approx$\,\textbf{84.8\,M} \\
    \bottomrule
  \end{tabular}
  \caption{Fuzuli model dimensions.}
  \label{tab:arch}
\end{table}

\subsection{Architectural choices in context}

Each component reflects current best-practice in dense decoder models
post-2023:
\begin{itemize}[leftmargin=*,nosep]
  \item \textbf{Decoder-only over encoder-decoder.} Following the GPT
    family~\citep{brown2020gpt3} and contemporary open-weight LMs
    (LLaMA~\citep{touvron2023llama2}). Encoder-decoder
    architectures~\citep{raffel2020t5} dominate paired-task settings;
    foundation-style next-token prediction has converged on decoder-only.
  \item \textbf{Grouped-Query Attention (GQA).} From
    \citet{ainslie2023gqa}: 4 KV heads against 8 query heads (a 2:1
    ratio), in the LLaMA-2 design idiom.
  \item \textbf{Rotary Position Embedding (RoPE).} From
    \citet{su2021roformer}.
  \item \textbf{SwiGLU MLP.} From \citet{shazeer2020glu}. The
    \texttt{mlp\_ratio} of 2.67 is the standard SwiGLU ratio that matches
    the parameter count of a GELU FFN at ratio 4.0.
  \item \textbf{RMSNorm.} From \citet{zhang2019rmsnorm}.
  \item \textbf{Untied input/output embeddings.} GPT-2 and many small
    models tie embedding and unembedding matrices for parameter efficiency.
    Fuzuli leaves them untied at the cost of an extra 15.4\,M parameters;
    the autoresearch loop tried and rejected tied embeddings.
  \item \textbf{No bias terms.} Following~\citep{touvron2023llama2}.
\end{itemize}

The training core itself is forked from Karpathy's
\texttt{nanochat}~\citep{karpathy2024nanochat}, stripped to single-file
form --- a deliberate choice that keeps the entire model definition readable
in one file, which matters for an autoresearch loop where an LLM-based
agent must comprehend and edit the training code.

\subsection{Training}

\begin{table}[t]
  \centering\small
  \begin{tabular}{ll}
    \toprule
    \textbf{Hyperparameter} & \textbf{Value} \\
    \midrule
    Total tokens seen & 60{,}000{,}000 ($\approx$ 7 epochs over 8\,M unique) \\
    Sequence length & 2048 \\
    Effective batch & $16 \times 4 = 64$ sequences \\
    Optimizer & AdamW, $\beta = (0.9, 0.95)$, $\text{wd}=0.1$ \\
    Learning rate & $6 \times 10^{-4}$ peak, cosine, 500 warmup steps \\
    Precision & bfloat16 \\
    Hardware & Single RTX 4090 \\
    Wall clock & $\approx$\,1--2 hours \\
    \bottomrule
  \end{tabular}
  \caption{Training hyperparameters for the promotion checkpoint.}
  \label{tab:train}
\end{table}

This is significantly under Chinchilla-optimal~\citep{hoffmann2022chinchilla}
--- for an 84\,M-parameter model the Chinchilla rule recommends $\approx$1.7
billion unique tokens. Section~\ref{sec:scaling} returns to this point.

\subsection{Evaluation at promotion scale}

The 84.8\,M-parameter promotion checkpoint was evaluated against the same
frozen heldout split used throughout the sprint search.

\begin{table}[t]
  \centering\small
  \begin{tabular}{lrrl}
    \toprule
    \textbf{Metric} & \textbf{Sprint baseline} & \textbf{Promotion} & $\boldsymbol{\Delta}$ \\
                    & \textbf{($\approx$30\,M)} & \textbf{(84.8\,M)} & \\
    \midrule
    Heldout PPL (bits/byte)        & 2.199 & \textbf{1.929} & $-0.270$ (improved) \\
    Ottoman lexicon score          & 2.677 & 9.873          & $+7.20$ (regressed) \\
    Modern-Turkish flatness        & 0.0037 & 0.0110        & $+0.007$ (regressed) \\
    Smoke fail rate                & 0.40 (2/5) & 0.80 (4/5) & $+0.40$ (regressed) \\
    \bottomrule
  \end{tabular}
  \caption{Sprint search baseline vs.\ promotion checkpoint on the frozen
    four-metric harness. The larger model improves perplexity by 0.27
    bits/byte but regresses on every generation-time metric --- the
    canonical data-ceiling pattern.}
  \label{tab:promotion-eval}
\end{table}

The 84.8\,M model \textbf{improves heldout perplexity by 0.27 bits/byte}
--- meaningful evidence that the larger architecture has captured more of
the corpus's statistical structure. But it regresses on every
generation-time metric.

This divergence is the v0.1 data-ceiling pattern restated.
Heldout perplexity is a \emph{closed-form} metric: it scores the model's
likelihood on text the model never sees during generation. The other
metrics depend on what the model produces \emph{under autoregressive
sampling}. A larger model has more capacity to memorize the corpus (lower
PPL) but, without enough data to ground its generation, correspondingly
more freedom to drift under sampling. The smaller sprint model is more
conservative simply because it has less capacity to express creative drift
with. Expanding the corpus to v0.5 / v1.0 targets is the only intervention
that recovers the generation metrics; no recipe tweak at this data scale
will.

% =============================================================================
\section{Autoresearch methodology}\label{sec:autoresearch}
% =============================================================================

The recipe was not hand-tuned. Instead, an LLM-based research agent was
given write access to a constrained surface --- \texttt{train/configs/}
\texttt{sprint.yaml}, \texttt{train/train.py}, \texttt{train/arch.py},
\texttt{data/cleaning\_rules.yaml} --- and instructed to propose one focused
change per cycle, run a 5-minute sprint training, observe the results, and
iterate. Around this agent sits a \textbf{factory}: a deterministic
orchestrator that enforces protected-path immutability, scans patches for
forbidden architectural patterns, runs the training, evaluates against four
metrics, and applies a strict-no-trades accept/reject policy. Accepted
experiments are merged into \texttt{main}; rejected experiments have their
branches deleted but their full record preserved in a SQLite ledger.

\begin{figure}[t]
  \centering
  \includegraphics[width=\textwidth]{\graphicsdir/system_architecture.png}
  \caption{Autoresearch system: the LLM agent proposes patches; the factory
    enforces protected paths, runs the training subprocess, evaluates, and
    applies the strict-no-trades policy; SQLite serves as cross-session
    memory.}
  \label{fig:sysarch}
\end{figure}

\subsection{Design choices in context}\label{sec:design}

Most of Fuzuli's individual components are off-the-shelf. The autoresearch
loop itself is part of a growing literature on LLM-driven scientific search:
FunSearch~\citep{romeraparedes2024funsearch} demonstrated that an LLM in a
closed evaluation loop could discover novel mathematical programs by
mutating a designated function while a fixed evaluator gates accept/reject;
AlphaEvolve~\citep{novikov2025alphaevolve} generalized this to algorithmic
discovery on a wider problem class; Voyager~\citep{wang2023voyager} and
Reflexion~\citep{shinn2023reflexion} established LLM-agent patterns for
persistent memory and self-correction across episodes. What this project
contributes is not the autoresearch primitive itself but a specific
combination of design decisions tailored to the sub-Chinchilla single-GPU
language-modeling setting:
\begin{itemize}[leftmargin=*,nosep]
  \item \textbf{Hard tier-based mutability.} The evaluation harness
    (\texttt{eval/}), the tokenizer artifact, the orchestrator
    (\texttt{factory/}), the CLI, and a handful of policy-defining files are
    \emph{immutable to the agent}. Any patch that modifies them is rejected
    at pre-flight by glob-based protected-path checking.
  \item \textbf{Strict-no-trades policy combiner.} The accept rule is
    purely Pareto: improve at least one metric beyond an improvement
    threshold, regress no metric beyond a regression tolerance. There is no
    scalar fitness function.
  \item \textbf{SQLite ledger as agent memory.} Every experiment ever run
    --- accept, reject\_eval, reject\_smoke, scores, deltas, reject reasons,
    branch SHAs --- is preserved in \texttt{experiments.sqlite}. A fresh
    agent session reads the ledger at startup and recovers full memory of
    what has been tried.
  \item \textbf{Forbidden-pattern scanning before GPU.} The agent cannot
    import \texttt{torch.distributed}, define MoE, Mamba, S4, Hyena, or any
    encoder module. The factory regex-scans every patch before allocating a
    GPU; the wall-clock cost of bad patches is bounded to under 30 seconds.
\end{itemize}

The four metrics:
\begin{itemize}[leftmargin=*,nosep]
  \item \texttt{heldout\_ppl}: bits-per-byte on a frozen heldout split.
  \item \texttt{lexicon\_score}: $-\log$ of the fraction of generated
    tokens present in an Ottoman lexicon.
  \item \texttt{flatness}: modern-Turkish loanword rate in generations.
  \item \texttt{smoke}: fraction of fixed prompts failing deterministic
    pass/fail rules.
\end{itemize}

% =============================================================================
\section{Results from the first 24 hours}
% =============================================================================

Over its first 24 hours of operation, the agent ran \textbf{356
experiments}: 6 accepts, 343 evaluation rejects, 7 smoke crashes (patches
that broke training before the eval phase), and zero infrastructure errors.

\begin{figure}[t]
  \centering
  \includegraphics[width=\textwidth]{\graphicsdir/score_progression.png}
  \caption{All four evaluators across 356 experiments. Stars mark accepted
    baselines; the dark step lines trace the current baseline's score on
    each metric. Most improvement happens in the first ${\sim}14$ experiments;
    the metric trajectories then flatten.}
  \label{fig:scores}
\end{figure}

The baseline trajectory makes the story visible: the productive period is
the first ${\sim}14$ experiments, during which lexicon score crashes from
10.05 to 2.85 --- a 71\% reduction driven primarily by a sequence-length
change ($1024 \to 512$) and a weight-decay tightening ($0.1 \to 0.2$).
Heldout perplexity barely moves throughout, suggesting the model converges
to roughly the same data-distribution likelihood under any reasonable
recipe.

\begin{figure}[t]
  \centering
  \includegraphics[width=\textwidth]{\graphicsdir/accept_rate.png}
  \caption{Rolling 20-experiment accept rate. Peaks above 30\% in the early
    cycles; collapses toward zero by experiment $\sim$50; after experiment
    $\sim$200, no further accepted experiments at all.}
  \label{fig:acceptrate}
\end{figure}

\begin{figure}[t]
  \centering
  \includegraphics[width=0.95\textwidth]{\graphicsdir/reject_reasons.png}
  \caption{Breakdown of rejection reasons across all non-accepts. PPL
    regressions dominate; smoke crashes are seven cases where a patch broke
    training before the eval phase.}
  \label{fig:rejects}
\end{figure}

% =============================================================================
\section{What the agent tried}
% =============================================================================

All 356 experiments were run on the \textbf{sprint configuration} --- a
smaller proxy model ($n_\text{layers}=6$, $n_\text{embd}=384$,
$n_\text{kv\_heads}=2$, $\approx$30\,M parameters, 25\,M training tokens,
$\approx$76\,s wall clock) used to discover hyperparameter recipes that are
then applied to the larger $\approx$84\,M-parameter promotion model in
Section~\ref{sec:architecture}. This separation is by design: sprints are
short enough for hundreds of cycles per day; promotions are 1--2 hours
each. The autoresearch loop iterates only over sprints.

Over the session, the agent traversed essentially the entire standard
recipe-search playbook: model width, layer count, KV-head count, peak
learning rate, learning-rate schedule (cosine vs constant), warmup-step
count, weight decay, initialization standard deviation, embedding tying,
dropout, sequence length, total training tokens, era data mix, random
seed, and optimizer choice (AdamW vs Lion~\citep{chen2023lion}). Most
productive accepts cluster around the early experiments.

% =============================================================================
\section{Diagnosis: data ceiling, not recipe ceiling}
% =============================================================================

The most informative single experiment is the one the agent ran toward the
end, when it began varying only the random seed.

\begin{figure}[t]
  \centering
  \includegraphics[width=0.85\textwidth]{\graphicsdir/seed_variance.png}
  \caption{Same recipe, five different \texttt{torch.manual\_seed} values.
    Lexicon score swings from 2.85 (seed 42) to 9.96 (seed 1) --- a
    ${\sim}3.5\times$ swing produced by initialization noise alone.}
  \label{fig:seedvar}
\end{figure}

Five seeds, identical configuration. The lexicon score under seed 42 is
2.85; under seed 1 it is 9.96. That variance is enormous relative to the
0.05 improvement threshold the policy requires for acceptance. A run that
happens to land in a ``good basin'' looks dramatically better than a run
that lands in a ``bad basin,'' and no recipe lever consistently moves the
model into the good basins. With only five seeds the evidence is
suggestive rather than conclusive, but the pattern is consistent with the
textbook signature of a \emph{data-starved training regime}.

\subsection{Where Fuzuli sits in the scaling landscape}\label{sec:scaling}

\begin{figure}[t]
  \centering
  \includegraphics[width=\textwidth]{\graphicsdir/scaling_landscape.png}
  \caption{Parameter count vs.\ unique training tokens (log-log) for Fuzuli
    v0.1 and a panel of public reference points. The dashed line is the
    Chinchilla~\citep{hoffmann2022chinchilla} compute-optimal ratio of 20
    tokens/parameter; the shaded regions mark over- and under-trained
    regimes. Fuzuli is firmly in the low-resource band alongside
    TinyStories~\citep{eldan2023tinystories}. \emph{Note:} BERTurk is an
    MLM-pretrained encoder; the Chinchilla curve was derived on CLM training,
    so the BERTurk position is approximate.}
  \label{fig:scaling}
\end{figure}

Chinchilla established that for compute-optimal training, the number of
unique training tokens should scale roughly as $20\times$ the number of
parameters. Fuzuli has \textbf{8\,M unique training tokens for 84\,M
parameters --- a unique-data-to-parameter ratio of $\approx$0.095}, against
the Chinchilla optimum of 20. That is roughly \textbf{$\mathbf{200\times}$
under} the compute-optimal frontier. Counting training-token \emph{passes}
across the $\sim$7 epochs raises the seen-token ratio to $\sim$0.71, but
Chinchilla scaling is properly defined on unique data;
\citet{muennighoff2023scaling} examine the data-constrained regime where
multiple epochs partly but not fully recover the marginal value of
additional unique data.

The opposite end of the landscape is illustrative. LLaMA-2-7B was trained
on 2 trillion tokens for 7 billion parameters~\citep{touvron2023llama2} ---
a ratio of $\sim$286, fully $14\times$ \emph{over} the Chinchilla optimum.
The modern best practice for production LMs is intentional overtraining;
Fuzuli takes the opposite trade because there is no deployment-scale
inference and there is no 2T-token Ottoman corpus.

The corollary is that further sprint iteration on the current corpus will
not produce meaningful gains. The recipe space is saturated; the data is
the binding constraint.

% =============================================================================
\section{Honest limitations of v0.1}
% =============================================================================

The model \textbf{will}:
\begin{itemize}[leftmargin=*,nosep]
  \item generate plausible Ottoman-Turkish continuations for short prompts;
  \item score perplexity on novel Ottoman text reliably enough for
    OCR-confidence and era-proximity scoring;
  \item provide useful Ottoman embeddings if a mean-pooling head is
    attached;
  \item demonstrate that a from-scratch Latinized Ottoman LM at this size
    is feasible.
\end{itemize}
The model \textbf{will not}:
\begin{itemize}[leftmargin=*,nosep]
  \item produce coherent long-form generation; expect memorized 3--5-token
    windows from the training corpus;
  \item handle out-of-distribution text well --- Arabic-script Ottoman,
    modern Turkish, conversational registers all degrade sharply;
  \item exhibit world knowledge beyond what is implicit in the corpus;
  \item replace human judgment on any historical-translation or
    paleographic task.
\end{itemize}
This is a research artifact. Treat it as one.

% =============================================================================
\section{Roadmap}
% =============================================================================

\begin{figure}[t]
  \centering
  \includegraphics[width=0.85\textwidth]{\graphicsdir/roadmap.png}
  \caption{Corpus tokens and model size targets for v0.5 (M2) and v1.0 (M3).
    Each version scales together with the corpus.}
  \label{fig:roadmap}
\end{figure}

Future versions of Fuzuli scale together with the corpus. v0.5 targets
$\sim$30\,M unique tokens at $\sim$120\,M parameters; v1.0 targets
$\sim$50\,M unique tokens at $\sim$150\,M parameters. Corpus growth is the
main bottleneck. Sources under consideration include the remaining un-OCR'd
late-Ottoman periodicals (\emph{İkdam}, \emph{Tasvir-i Efkâr}, \emph{Mizan},
\emph{Tercümân-ı Ahvâl}), Latinized Ottoman primary-source passages
embedded in Turkish-studies dissertations available through the YÖK Tez
archive, and academic encyclopedia entries from TDV İslâm Ansiklopedisi.

A parallel encoder track (BERT-style~\citep{devlin2019bert} masked
language-model pretraining on the same corpus) is also planned. Encoder
objectives are substantially more sample-efficient than decoder objectives;
the existing 8\,M tokens are comfortable for a 50--80\,M-parameter encoder.

% =============================================================================
\section*{License}
% =============================================================================

\begin{itemize}[leftmargin=*,nosep]
  \item \textbf{Code}: Apache-2.0
  \item \textbf{Model weights}: Apache-2.0
  \item \textbf{Dataset}: CC-BY-4.0
  \item \textbf{This article and its diagrams}: CC-BY-4.0
\end{itemize}

This work is released without external attribution claims. Source repository:
\url{https://github.com/fbkaragoz/asena-project}.

% =============================================================================
\bibliographystyle{plainnat}
\begin{thebibliography}{99}

\bibitem[Ainslie et al.(2023)]{ainslie2023gqa}
Ainslie, J., Lee-Thorp, J., de Jong, M., Zemlyanskiy, Y., Lebrón, F.,
\& Sanghai, S. (2023).
\newblock GQA: Training Generalized Multi-Query Transformer Models from Multi-Head Checkpoints.
\newblock \emph{EMNLP 2023}.

\bibitem[Brown et al.(2020)]{brown2020gpt3}
Brown, T. B. \emph{et al.} (2020).
\newblock Language Models are Few-Shot Learners.
\newblock \emph{NeurIPS 2020}.

\bibitem[Chen et al.(2023)]{chen2023lion}
Chen, X. \emph{et al.} (2023).
\newblock Symbolic Discovery of Optimization Algorithms.
\newblock \emph{arXiv:2302.06675}.

\bibitem[Devlin et al.(2019)]{devlin2019bert}
Devlin, J., Chang, M.-W., Lee, K., \& Toutanova, K. (2019).
\newblock BERT: Pre-training of Deep Bidirectional Transformers for Language Understanding.
\newblock \emph{NAACL 2019}.

\bibitem[Eldan \& Li(2023)]{eldan2023tinystories}
Eldan, R., \& Li, Y. (2023).
\newblock TinyStories: How Small Can Language Models Be and Still Speak Coherent English?
\newblock \emph{arXiv:2305.07759}.

\bibitem[Hoffmann et al.(2022)]{hoffmann2022chinchilla}
Hoffmann, J. \emph{et al.} (2022).
\newblock Training Compute-Optimal Large Language Models.
\newblock \emph{arXiv:2203.15556} (Chinchilla).

\bibitem[Karagöz et al.(2024)]{karagoz2024towards}
Karagöz, F., Doğan, B., \& Özateş, Ş.~B. (2024).
\newblock Towards a Clean Text Corpus for Ottoman Turkish.
\newblock \emph{Proceedings of the First Workshop on Natural Language Processing for Turkic Languages (SIGTURK 2024)}, 62--70.

\bibitem[Karpathy(2024)]{karpathy2024nanochat}
Karpathy, A. (2024--).
\newblock nanoGPT / nanochat training core.
\newblock \url{https://github.com/karpathy/nanoGPT}.

\bibitem[Muennighoff et al.(2023)]{muennighoff2023scaling}
Muennighoff, N. \emph{et al.} (2023).
\newblock Scaling Data-Constrained Language Models.
\newblock \emph{NeurIPS 2023}.

\bibitem[Novikov et al.(2025)]{novikov2025alphaevolve}
Novikov, A. \emph{et al.} (2025).
\newblock AlphaEvolve: A coding agent for scientific and algorithmic discovery.
\newblock \emph{DeepMind / arXiv}.

\bibitem[Özateş et al.(2025)]{ozates2025foundations}
Özateş, Ş.~B., Tıraş, T.~E., Adak, E.~E., Doğan, B., Karagöz, F.~B.,
Genç, E.~E., \& Taşdemir, E.~F.~B. (2025).
\newblock Building Foundations for Natural Language Processing of Historical Turkish: Resources and Models.
\newblock \emph{arXiv:2501.04828}.

\bibitem[Raffel et al.(2020)]{raffel2020t5}
Raffel, C. \emph{et al.} (2020).
\newblock Exploring the Limits of Transfer Learning with a Unified Text-to-Text Transformer (T5).
\newblock \emph{JMLR 21(140)}.

\bibitem[Romera-Paredes et al.(2024)]{romeraparedes2024funsearch}
Romera-Paredes, B. \emph{et al.} (2024).
\newblock Mathematical discoveries from program search with large language models (FunSearch).
\newblock \emph{Nature 625}, 468--475.

\bibitem[Schweter(2020)]{schweter2020berturk}
Schweter, S. (2020).
\newblock BERTurk — BERT models for Turkish.
\newblock \emph{Zenodo}. doi:10.5281/zenodo.3770924.

\bibitem[Sennrich et al.(2016)]{sennrich2016bpe}
Sennrich, R., Haddow, B., \& Birch, A. (2016).
\newblock Neural Machine Translation of Rare Words with Subword Units.
\newblock \emph{ACL 2016}.

\bibitem[Shazeer(2020)]{shazeer2020glu}
Shazeer, N. (2020).
\newblock GLU Variants Improve Transformer.
\newblock \emph{arXiv:2002.05202}.

\bibitem[Shinn et al.(2023)]{shinn2023reflexion}
Shinn, N., Cassano, F., Berman, E., Gopinath, A., Narasimhan, K., \& Yao, S. (2023).
\newblock Reflexion: Language Agents with Verbal Reinforcement Learning.
\newblock \emph{NeurIPS 2023}.

\bibitem[Su et al.(2021)]{su2021roformer}
Su, J., Lu, Y., Pan, S., Murtadha, A., Wen, B., \& Liu, Y. (2021).
\newblock RoFormer: Enhanced Transformer with Rotary Position Embedding.
\newblock \emph{arXiv:2104.09864} (RoPE).

\bibitem[Suárez et al.(2019)]{suarez2019oscar}
Suárez, P.~J.~O., Romary, L., \& Sagot, B. (2019).
\newblock Asynchronous Pipeline for Processing Huge Corpora on Medium to Low Resource Infrastructures (OSCAR).
\newblock \emph{Workshop on Challenges in the Management of Large Corpora}.

\bibitem[Touvron et al.(2023)]{touvron2023llama2}
Touvron, H. \emph{et al.} (2023).
\newblock LLaMA 2: Open Foundation and Fine-Tuned Chat Models.
\newblock \emph{arXiv:2307.09288}.

\bibitem[Vaswani et al.(2017)]{vaswani2017attention}
Vaswani, A. \emph{et al.} (2017).
\newblock Attention Is All You Need.
\newblock \emph{NeurIPS 2017}.

\bibitem[Wang et al.(2023)]{wang2023voyager}
Wang, G. \emph{et al.} (2023).
\newblock Voyager: An Open-Ended Embodied Agent with Large Language Models.
\newblock \emph{arXiv:2305.16291}.

\bibitem[Zhang \& Sennrich(2019)]{zhang2019rmsnorm}
Zhang, B., \& Sennrich, R. (2019).
\newblock Root Mean Square Layer Normalization.
\newblock \emph{NeurIPS 2019}.

\end{thebibliography}

\end{document}
